\chapter{MPS/DMRG support in the release}
\label{ch:mps-support}

\section{Purpose and scope}

The release tree does not contain a tensor-network DMRG solver, and it
does not contain the material-diffusion tools of the development
branch.  What it contains is one graduated DMRG-line utility,
\texttt{core.mps.pin\_degenerate\_columns}, which resolves a
reproducibility problem that appears before any DMRG calculation can be
run: pyscf RHF orbitals are unique only up to arbitrary orthogonal
rotations inside exactly degenerate irreps, so two processes can return
different representatives of the same degenerate pair.  A
column-referenced active-space window, such as ``columns 6 and 7'', then
means different orbitals in different runs.

The utility applies to any column-referenced orbital window on a
canonical SCF or CASSCF orbital set.  It does not select an active
space, rank windows, or run a DMRG calculation; it only fixes the
representative of a degenerate group.  It is NumPy-only and has no JAX
or pyscf dependency.

\section{Theory and derivation}
\label{sec:pin-theory}

\subsection{Degenerate subspaces and gauge freedom}

Let $F$ and $S$ be the Fock and overlap matrices in the AO basis, and
let $C$ be a set of orthonormal molecular orbitals,
\begin{equation}
  C^{T}SC = I,
  \qquad
  F C = S C\,\epsilon,
  \label{eq:pin-eigen}
\end{equation}
with $\epsilon=\mathrm{diag}(\epsilon_1,\dots,\epsilon_{n_{\mathrm{mo}}})$.
Suppose a group of columns $G=\{g_1,\dots,g_m\}$ has
\begin{equation}
  |\epsilon_i-\epsilon_j|<\tau_{\mathrm{deg}},
  \qquad i,j\in G,
  \label{eq:pin-group}
\end{equation}
with $\tau_{\mathrm{deg}}$ the \texttt{deg\_tol} parameter.  For any
orthogonal matrix $U\in O(m)$ the rotated set
\begin{equation}
  C'_G = C_G\,U
  \label{eq:pin-rotation}
\end{equation}
is an equally valid orthonormal eigenbasis of the same degenerate
subspace: the span is unchanged, the Fock eigenvalues are unchanged,
and the energy of any wavefunction built from the complete group is
unchanged.  Different SCF runs, different linear-algebra backends, or
different starting guesses can return different $U$.

If a column-referenced window cuts a degenerate group, the selected
columns are not a well-defined subspace.  On N$_2$/6-31G the RHF
spectrum carries the $\pi_u$ pair at columns $(5,6)$, the $\pi_g$ pair
at $(7,8)$, and two higher pairs.  A CAS(2,2) on the recorded start
window $\{6,7\}$ scattered over $0.02$~Ha across fresh builds although
the RHF energy matched to $3\times10^{-13}$, because the occupied
$\pi_u$ mixture decides which combination sits in the frozen core.

\subsection{The pinning map}

The utility picks a deterministic basis inside each degenerate group
by diagonalizing a caller-supplied Hermitian operator $X^{2}$ in the
group's span:
\begin{equation}
  \left(C_G^{T}X^{2}C_G\right)V = V\,\Lambda,
  \qquad
  \Lambda=\mathrm{diag}(\lambda_1,\dots,\lambda_m),
  \label{eq:pin-x2}
\end{equation}
and replaces
\begin{equation}
  C_G \leftarrow C_G\,V_{[:,\pi]},
  \label{eq:pin-replace}
\end{equation}
where $\pi$ orders the eigenvalues in descending order.  The
expectation values $\lambda_i$ are the diagonal $X^{2}$ values of the
new columns, and the columns are ordered by that value.  The phase of
each resulting column is fixed by making its largest-magnitude
coefficient positive:
\begin{equation}
  c_i \leftarrow \mathrm{sign}\!\left(c_{i,\mu(i)}\right)c_i,
  \qquad
  \mu(i)=\arg\max_\mu |c_{i\mu}|.
  \label{eq:pin-phase}
\end{equation}
Non-degenerate columns pass through unchanged except for this phase
convention.

The map has three properties that make it safe to insert into an
existing workflow:

\begin{itemize}
\item \emph{Span invariance.}  A window that contains a degenerate group
  whole is invariant under Eq.~\eqref{eq:pin-replace}: the span is
  preserved, so its CI energies do not move.
\item \emph{Determinism.}  A split group becomes deterministic.  Any
  single representative of a degenerate pair is symmetry-equivalent,
  but the process-to-process arbitrariness is removed.
\item \emph{Idempotence and gauge invariance.}  Pinning is idempotent
  and invariant to an input rotation applied inside a degenerate group:
  if $C_G$ is replaced by $C_GU$, the projected operator
  $U^{T}(C_G^{T}X^{2}C_G)U$ has the same eigenvectors transformed by
  $U$, and the final columns agree.
\end{itemize}

\subsection{Conditions and failure modes}

The construction is unique only when the restriction of $X^{2}$ to the
group has a non-degenerate spectrum.  If two eigenvalues of
$\Lambda$ coincide, the eigenvectors inside that sub-block are again
arbitrary, and the pinning cannot resolve the remaining gauge.  The
energy of a complete window is still invariant, but a column-split
window is not.  On a linear molecule the $x^{2}$ operator distinguishes
the two components of a $\pi$ pair, which is why it is the operator
used in the recorded calibration; a spherical operator would not.

The grouping is an energy window, not a symmetry classification.  A
\texttt{deg\_tol} that is too small splits a physical near-degeneracy
and pins only part of the group; one that is too large merges distinct
states.  The recorded test checks both sides of the tolerance
($2.0+0.5\times10^{-7}$ inside, $2.0+5.0\times10^{-7}$ outside the
default $10^{-7}$ window).

\section{Input and operation}
\label{sec:pin-input}

\begin{center}
\small
\begin{tabularx}{\textwidth}{>{\raggedright\arraybackslash}p{0.36\textwidth}X}
\toprule
Entry point & Returns and purpose \\
\midrule
\texttt{pin\_degenerate\_columns} &
  Arguments \texttt{(mo\_coeff, mo\_energy, x2\_ao, deg\_tol=1e-7)}.
  Returns \texttt{(C\_pinned, n\_rotated)}.  Every group of columns
  whose energies agree within \texttt{deg\_tol} is rotated into the
  eigenbasis of \texttt{x2\_ao}, ordered by descending expectation
  value; each column's phase is then fixed by making its
  largest-magnitude entry positive. \\
\bottomrule
\end{tabularx}
\end{center}

\texttt{mo\_coeff} has shape \texttt{(nao, nmo)} and is column-ordered;
\texttt{mo\_energy} has length \texttt{nmo}; \texttt{x2\_ao} is the
\texttt{(nao, nao)} matrix of the chosen operator in the same AO basis,
and \texttt{deg\_tol} is in the energy unit of \texttt{mo\_energy}
(Eh for SCF orbital energies).  \texttt{n\_rotated} counts the columns
inside degenerate groups that were re-rotated; the phase convention is
applied to all columns but is not counted.  The function returns a new
array and does not modify the input.

A minimal call is
\begin{lstlisting}[language=Python]
import numpy as np
from pyscf import gto, scf
from diffmrc.core.mps import pin_degenerate_columns

mol = gto.M(atom='N 0 0 0; N 0 0 1.10', basis='6-31g',
            charge=0, spin=0, verbose=0)
mf = scf.RHF(mol).run(verbose=0)

rx = mol.intor('int1e_r')[0]
C, n_rot = pin_degenerate_columns(np.asarray(mf.mo_coeff),
                                  np.asarray(mf.mo_energy), rx @ rx)
print('re-rotated columns', n_rot)
\end{lstlisting}

\section{Worked example}
\label{sec:pin-example}

\subsection{N\texorpdfstring{$_2$}{2}/6-31G: the recorded start window
\{6,7\}}

On N$_2$/6-31G at $R=1.10$~\AA{} the raw RHF energy is
$-108.867618373058$~Eh.  The call above returns \texttt{n\_rot}=8:
four degenerate pairs are re-rotated.  A CAS(2,2) on the pinned
$\{6,7\}$ window gives
\begin{equation}
  E_{\mathrm{CAS(2,2)}} = -108.8679136653~\text{Eh},
  \label{eq:pin-n2-window}
\end{equation}
stable across fresh builds to the printed digits.  Without pinning the
same window scattered over $0.02$~Ha.  The regression is the
\texttt{test\_n2\_6\_31g\_start\_window\_reference} case in
\texttt{tests/core/test\_mps.py}.

\subsection{N\texorpdfstring{$_2$}{2}/6-31G: the exhaustive-window
best}

The same pinned basis reproduces the exhaustive-window best of the
calibration:
\begin{equation}
  E_{\mathrm{elec}} = -132.4777614802~\text{Eh},
  \qquad
  E_{\mathrm{tot}} = -108.9053220847~\text{Eh},
  \label{eq:pin-n2-best}
\end{equation}
to $8\times10^{-11}$.  The two examples differ in the window, not in
the pinning: the first is the recorded start window, the second is the
best window found by the exhaustive scan.  The module docstring records
the calibration and the reproduction margins; the test pins the first
window, and the second value is the reference quoted in the same
graduation record.

\section{Pitfalls}
\label{sec:pin-pitfalls}

\begin{itemize}
\item \emph{A split group.}  Pinning a subset of a degenerate group
  preserves that subset's span but not the physical subspace.  The
  window must contain the group whole, or the calculation must use a
  window defined after pinning.
\item \emph{A degenerate \texttt{x2\_ao} restriction.}  If the chosen
  operator has equal expectation values inside the group, the
  eigenbasis is again arbitrary.  The $x^{2}$ operator works for the
  $\pi$ pairs of a linear molecule; a spherical operator does not.
\item \emph{The tolerance.}  \texttt{deg\_tol} is an energy window, not
  a symmetry test.  A value below the numerical splitting of a
  physical degeneracy leaves the group split; a value above a physical
  gap merges distinct states.
\item \emph{Non-Hermitian or non-Hermitianized \texttt{x2\_ao}.}  The
  eigenproblem in Eq.~\eqref{eq:pin-x2} assumes a Hermitian
  restriction; a non-symmetric input can return complex or unordered
  eigenvectors.
\item \emph{Wrong basis.}  \texttt{x2\_ao} must be in the same AO basis
  as \texttt{mo\_coeff}.  A Cartesian component from a different
  basis or a different unit system defines a different deterministic
  basis; the span invariance still holds, but the recorded reference
  values do not transfer.
\item \emph{It is not a DMRG solver.}  The function fixes orbital
  representatives for column-referenced windows; the DMRG candidate
  space and the material-diffusion tools mentioned in the development
  notes are not part of this release tree.
\end{itemize}

\section{Validation and references}
\label{sec:pin-validation}

\texttt{tests/core/test\_mps.py} pins the synthetic contracts (pair
axes, span, phases, idempotence, invariance to an input rotation,
tolerance behaviour, and the no-group case) and the N$_2$/6-31G
molecule regression: RHF $-108.867618373058$~Eh, \texttt{n\_rot}$=8$,
pinned CAS(2,2) $-108.8679136653$~Eh.  The test needs pyscf for the
molecular case and NumPy only for the synthetic cases.  The graduation
record is the module docstring of \texttt{core.mps}; the original
measurements were made in the T5 prototype driver
(\texttt{code/t416\_t5\_paper\_systems.py} in the development tree) and
are summarized in the T5 report, neither of which is part of this
release tree.

The construction is the standard resolution of an orbital gauge
freedom inside a degenerate subspace; the pyscf orbital conventions and
the Fock eigenproblem are documented in the pyscf \texttt{scf} and
\texttt{mcscf} documentation.  The $x^{2}$-based pinning was chosen for
the linear-molecule $\pi$ pairs because the $x^{2}$ restriction has a
real eigenvalue gap inside the pair; the calibration and the
reproduction margins are the record cited above.
